\section{问题一：确定性线性规划模型}

\subsection{模型建立}

问题一中电价、负载与光伏预测全天已知，储能首尾电量相等，是一个确定性优化问题。以全天购电费用最小为目标，决策变量为 144 个区间的购电量 $x_t$、充电量 $c_t$、放电量 $d_t$、储电量 $s_t$ 与弃光量 $w_t$。

\begin{equation}
\min_{x,c,d,s,w}\quad \sum_{t=0}^{N-1} p_t x_t
\label{eq:p1obj}
\end{equation}

约束条件如下。

\textbf{功率平衡}（等式，不允许售电）：
\begin{equation}
x_t+PV_t+d_t-c_t-w_t=L_t,\qquad t=0,\dots,N-1 .
\label{eq:p1balance}
\end{equation}

\textbf{储能递推与容量}：
\begin{equation}
s_t=s_{t-1}+\eta_c c_t-\frac{d_t}{\eta_d},\quad
s_{-1}=s_{N-1}=S_0=6000,\quad
S^{\min}\le s_t\le S^{\max}.
\label{eq:p1soc}
\end{equation}

\textbf{充放电功率与购电边界}：
\begin{equation}
0\le c_t\le P^{\max}\Delta t,\qquad
0\le d_t\le P^{\max}\Delta t,\qquad
x_t\ge0,\qquad w_t\ge0 .
\label{eq:p1bounds}
\end{equation}

式~\eqref{eq:p1balance}--\eqref{eq:p1bounds} 全部是线性等式或线性不等式，式~\eqref{eq:p1obj} 是线性目标，因此这是一个标准线性规划\cite{boyd2004}，采用 HiGHS 单纯形法求解\cite{virtanen2020}。变量共 $5\times144=720$ 个，等式约束 $2\times144+1=289$ 条。

\subsection{求解结果}

求解得到全天购电量 $57580.98\,\mathrm{kWh}$、全天购电费 $33416.53$ 元，储能 0:00 与 24:00 的储电量均为 $6000.00\,\mathrm{kWh}$，全天弃光为 $0$。表 2、表 3 给出论文要求的展示口径。

\threelinetable{表 2\quad 指定时间段购电量与全天合计}{lrlr}
{时间段 & 购电量 / kWh & 时间段 & 购电量 / kWh}
{10:00-10:10 & 0.0000 & 16:00-16:10 & 445.4317 \\
 12:00-12:10 & 480.4124 & 18:00-18:10 & 0.0000 \\
 14:00-14:10 & 0.0000 & 20:00-20:10 & 0.0000 \\
 \multicolumn{3}{l}{全天购电量 / kWh} & 57580.98 \\
 \multicolumn{3}{l}{全天购电费 / 元} & 33416.53}

\threelinetable{表 3\quad 储能设备分时段充放电量与首尾储电量}{lrr}
{时间段 & 充电量 / kWh & 放电量 / kWh}
{0:00-4:00   & 4500.0000 & 0.0000 \\
 4:00-8:00   & 833.3333  & 6956.9814 \\
 8:00-12:00  & 4448.6339 & 1702.9970 \\
 12:00-16:00 & 5274.7881 & 91.1014 \\
 16:00-20:00 & 0.0000    & 6740.1319 \\
 20:00-24:00 & 5333.3333 & 2859.8681 \\
 \multicolumn{2}{l}{0:00 储电量 / kWh} & 6000.00 \\
 \multicolumn{2}{l}{24:00 储电量 / kWh} & 6000.00}

\begin{figure}[H]
  \centering
  \includegraphics[width=0.94\textwidth]{p1_dispatch.pdf}
  \caption{问题一典型日的电价、计划购电量、储能充放电量与储电量}
  \label{fig:p1}
\end{figure}

图~\ref{fig:p1} 显示调度策略呈现清晰的“低价充、高价放”结构，储电量轨迹可以分为四段。

\textbf{夜间低价充电（0:00--5:30）。}此时电价处于 $0.42\sim0.43$ 元/$\mathrm{kWh}$ 的当日低位，储能以接近上限的功率充电，储电量从 $6000\,\mathrm{kWh}$ 升至 $10050\,\mathrm{kWh}$ 并于 5:30 触及 $10800\,\mathrm{kWh}$ 的上限。

\textbf{早高峰放电（5:30--10:00）。}光伏出力尚小，小区负载在 9:00 附近达到全天最大值 $993.16\,\mathrm{kWh}$/区间（对应 $5959\,\mathrm{kW}$），储能放电顶负荷。储电量从上限回落到 10:00 的 $2279\,\mathrm{kWh}$，其间 5:40--8:00 等 61 个区间购电量降为 $0$。

\textbf{白天光伏富余充电（10:00--16:00）。}全天有 30 个区间光伏出力超过小区负载，富余总量 $6248.0\,\mathrm{kWh}$，全部被储能吸收，没有弃光。储电量从 $2279\,\mathrm{kWh}$ 回升并在 16:00 重新触顶 $10800\,\mathrm{kWh}$，16:00--18:00 维持满电待发。

\textbf{晚高峰放电与回充（18:00--24:00）。}电价进入 $1.3$ 元/$\mathrm{kWh}$ 以上的峰段，储能集中放电，储电量在 20:30 触及 $1200\,\mathrm{kWh}$ 的下限；21:00 之后电价回落，储能再次以低价电补回电量，24:00 精确回到 $6000\,\mathrm{kWh}$。

全天弃光为 $0$，说明储能容量足以吸收全部光伏富余。论文要求的六个展示区间中，10:00--10:10、14:00--14:10、18:00--18:10、20:00--20:10 四个区间购电量均为 $0$，完全由光伏与储能供电。
